On considère la fonction $f$ définie sur $\R$ par :
\[f(x)=(4 x-4) \e^{-0,5 x}+5\]
On note $\mathcal{C}_{f}$ sa courbe représentative dans un repère orthonormé.

On admet que $f$ est dérivable sur $\R$ et on note $f'$ sa fonction dérivée.

\begin{enumerate}
	\item Montrer que, pour tout $x \in \R$,\quad $f'(x)=(-2 x+6) \e^{-0,5 x}$.

	\item Étudier le signe de $f'(x)$ sur $\R$ puis en déduire les variations de la fonction $f$ sur $\R$.

	\item La courbe $\mathcal{C}_{f}$ admet-elle des points pour lesquels la tangente est horizontale ? Si oui, on précisera les coordonnées exactes de ces éventuels points.
\end{enumerate}
